% chapters/02_literature_review.tex - 第2章 文献综述

\chapter{文献综述}

\section{反馈在抗阻训练中的作用与类型划分}

抗阻训练的适应效果并非由外部负荷单向决定，而是由训练者能否依据训练过程中获得的反馈信息对负荷进行动态调节所共同塑造。多模态AI辅助训练系统之所以能够在近年来的健身领域受到广泛关注，其核心价值也正在于通过通用智能设备为训练者提供更精准、更及时的反馈信息，从而降低传统专业监控设备的应用门槛。因此，理解此类AI辅助训练的科学价值，首先需要理解反馈本身在抗阻训练中的作用机制、类型划分及其代表方法的方法学内涵。本节围绕反馈方式这一主线，从训练适应的反馈依赖机制出发，依次论述反馈的生理学与行为学作用机制、反馈的类型划分、两类反馈的代表方法及其相互关系，为后续章节从学术史与证据梳理层面的深入评述奠定统一的概念基础。

\subsection{抗阻训练适应对反馈的依赖}

抗阻训练诱发神经肌肉适应的前提是训练刺激与训练者当日功能状态之间的有效匹配。\citep{Kraemer2004}在其经典综述中系统论证了训练强度、训练量与恢复安排三者之间的交互作用，指出训练适应的方向与幅度不仅取决于负荷的绝对大小，更取决于负荷是否落在与训练者当日神经肌肉能力相匹配的有效刺激区间内。这一论断的生理学基础在于，个体的最大力量表现在不同训练日之间存在显著的日内波动。\citep{Zourdos2016}的研究表明，即使是训练经验较为丰富的力量训练者，其不同训练日的1RM表现也可能出现10\%至18\%的波动，波动来源涵盖睡眠质量、生活压力、营养状况以及既往训练负荷的残余疲劳等多重因素。这种日内波动意味着，若训练者始终依据训练前预设的固定百分比处方执行训练，实际相对强度将不可避免地偏离目标区间。

要解决固定处方与日内波动之间的矛盾，训练者必须在训练过程中获取关于当前负荷是否适宜的即时信息，并据此对后续训练的负荷、组数或重复次数做出动态调整。\citep{Greig2020}将这一过程定义为自主调节训练（autoregulated resistance training），即"基于对个体表现的客观测量或其对执行能力的主观感知来调整训练的过程"。这一定义揭示了自主调节训练的两个核心特征。第一，训练负荷根据训练者当日状态实时调整，而非机械执行预设处方。第二，调节依据既可以是客观测量，也可以是主观感知，两者在逻辑上处于并列地位。反馈信息因此并非训练的附属品，而是连接外部训练刺激与内部神经肌肉适应的核心中介变量。

\subsection{反馈的作用机制与两种类型}

反馈之所以能够改善训练适应，其作用机制可从生理学与行为学两个层面加以理解。在生理学层面，反馈信息的核心功能在于帮助训练者识别当前神经肌肉状态，实现负荷强度与疲劳程度的精准匹配。\citep{GonzalezBadillo2011}的研究证实，杠铃向心速度的下降与神经肌肉疲劳积累和代谢应激水平之间存在稳定的剂量-效应关系，这意味着速度反馈能够为训练者提供关于疲劳累积程度的客观量化参照。在行为学层面，反馈信息通过调节训练者的动作意图与注意焦点影响神经肌肉输出。\citep{GonzalezBadillo2014}的研究表明，要求训练者在向心阶段以最大意图速度推动杠铃的神经指令本身即可驱动更高阈值运动单位的募集，而实时速度反馈能够强化这一意图指令的执行质量。\citep{Wulf2016}提出的OPTIMAL运动学习理论从更宏观的视角解释了反馈的行为学效应，该理论指出，增强期望（enhanced expectancies）、自主性支持（autonomy support）与外部注意焦点（external focus of attention）是优化运动表现与学习的三个核心因素。

基于反馈信息来源的根本差异，\citep{Greig2020}在对自主调节训练术语进行系统梳理时提出，抗阻训练中的反馈可划分为两大类型。第一类是客观测量反馈（objective feedback），其信息载体为外部设备采集的运动学或动力学参数，训练者依据这些参数判断负荷强度与疲劳状态。第二类是主观感知反馈（subjective feedback），其信息载体为训练者自身对用力感与剩余能力的内在评估，无需任何外部设备即可实施。这一划分反映了反馈信息获取的两种根本不同的技术路径。前者依赖外部工具将不可见的力学过程转化为可见的数字信号，后者依赖训练者将内在生理感受转化为可操作的判断。

\subsection{客观测量反馈的代表方法}

在客观测量反馈的诸多具体形式中，基于速度的训练（velocity-based training, VBT）是目前发展最为成熟、证据积累最为充分的方法。VBT的理论基础源于肌肉收缩中已被充分证实的生物力学规律，即力与速度之间的稳定负相关。\citep{GonzalezBadillo2010}通过对杠铃卧推的系统研究确立，给定相对负荷（\%1RM）与其对应向心速度之间存在高度可预测的线性关系。在此理论基础上，VBT发展出三类核心操作参数。第一是平均向心速度（mean concentric velocity, MCV），指杠铃在完整向心阶段的平均移动速度。第二是速度损失率（velocity loss, VL\%），指同一组内首次与末次重复速度之差占首次速度的百分比。第三是最大意图速度（maximum intentional velocity, MIV），指训练者在向心阶段以最大努力推动杠铃的神经指令。这三类参数共同界定了VBT作为一套完整方法学的操作内涵，构成了本研究客观反馈干预组的核心方法学基础。

\subsection{主观感知反馈的代表方法}

与VBT依赖外部设备采集客观数据形成对照，主观感知反馈路径通过训练者自身的内在评估实现负荷调节，其代表方法经历了从RPE到RIR的关键演进。\citep{Borg1982}最初提出的主观用力程度量表（rating of perceived exertion, RPE）通过整合呼吸困难、肌肉疲劳与主观努力程度等多维信号对运动强度进行整体量化。\citep{Zourdos2016}针对抗阻训练的动作特性开发了基于储备重复次数（repetitions in reserve, RIR）的RPE量表，将抽象的用力感转化为一个具有直接操作价值的判断，即在维持标准动作技术的前提下本组结束时尚能完成的极限重复次数。RIR相较于传统RPE的方法学优势体现在操作直接性、与动作力学的直接对应性以及与客观参数的可桥接性三个方面。

\section{客观反馈路径的方法学演进与实证进展}

客观反馈路径的发展跨越近九十年，其中有四个关键节点：力-速度关系的理论奠基、负荷-速度关系在训练场景中的验证、速度损失指标的方法学发展，以及计算机视觉技术的应用引入。本节梳理客观反馈路径从力-速度关系的实验室奠基到移动端计算机视觉技术推动的应用下沉的完整历史演进，并系统评述VBT相较于传统百分比训练的实证证据与技术实现现状。

\subsection{客观反馈路径的方法学演进}

\citep{Hill1938}通过青蛙缝匠肌实验首次确立肌肉收缩时张力与缩短速度之间的双曲线关系。此后\citep{GonzalezBadillo2010}将这一关系从实验室拓展到真实训练场景，使"通过速度反推强度"从理论可能变为操作现实。\citep{GonzalezBadillo2011}证实了速度衰减率与神经肌肉疲劳之间的剂量-效应关系，\citep{ParejaBlanco2017}进一步揭示不同速度损失阈值对应不同方向的适应。2020年后，深度学习姿态估计技术逐渐成熟（\citep{Bazarevsky2020}），使基于智能手机摄像头的VBT方案在原理上成为可能。

\subsection{VBT相较于传统百分比训练的适应性证据}

2025年发表的伞形综述确认，VBT相较于PBT在最大力量、爆发力与冲刺表现上的适应性优势已具备较为一致的高等级证据支持。\citep{Wlodarczyk2021}纳入15项随机对照试验，结果提示采用10\%至20\%速度损失阈值的VBT方案在四个维度上均具备优于PBT的适应潜力。\citep{Held2022}通过网络Meta分析进一步证实，低到中等速度损失阈值的VBT方案在爆发力与最大力量适应上均优于PBT。\citep{Luo2026}作为该领域最新的Meta分析，纳入17项随机对照试验与348名受试者，结果显示VBT在跳跃表现（SMD=0.27, $p$$<$0.05）与变向能力（SMD=0.45, $p$$<$0.01）上显著优于PBT，但在最大力量与直线冲刺上未达显著差异。

然而，上述所有证据的对照组均为传统PBT。VBT相较于PBT的适应性优势，可能部分来自VBT本身的方法学优势，部分则来自PBT作为参照方案的固有局限。这引出一个值得追问的问题：当对照组从PBT更换为同样具备科学基础的规范化RIR训练时，VBT是否仍能维持其已确立的适应性优势？

\subsection{客观反馈的技术实现与效度证据}

VBT的实际应用依赖速度测量工具的技术可行性与测量精度。当前客观反馈路径的技术实现可分为三类：基于物理位移测量的线性位置传感器（LPT）、基于惯性测量单元的可穿戴加速度计，以及基于图像识别的计算机视觉（CV）方案。三类方案在测量原理、精度水平、成本门槛与适用场景上存在明显差异。

线性位置传感器（LPT）是VBT领域应用历史最长、效度验证最充分的测量设备。\citep{Weakley2021}在涵盖22种商用VBT设备的系统综述中确认，GymAware在平均速度与峰值速度测量上的效度最高（相关系数$r$=0.90至1.00，SEE=0.01至0.08 m/s）。但单台商用LPT设备价格通常在2000美元以上，使LPT的应用长期集中于职业运动队、高水平训练中心和科研实验室。

可穿戴加速度计在便携性与成本上具备优势，但效度面临严峻挑战。\citep{Banyard2017}将PUSH Band与GymAware及实验室标准设备进行三方对比，结果显示GymAware对所有效标变量高度有效，而PUSH Band仅在峰值力估计上达到较高精度，在平均速度与峰值速度测量上精度不足。\citep{PerezCastilla2022}在自由重量与史密斯机深蹲中对PUSH Band 2.0进行了更系统的评估，研究结论明确指出PUSH Band 2.0是"可靠但无效"的可穿戴技术。

计算机视觉（CV）方案是近年来发展最快的技术路径。\citep{Renner2024}对三款主流商用VBT应用的并发效度评估揭示了CV方案的三类典型问题：漏检问题（在总共589次深蹲重复中，Metric VBT漏检52次（8.8\%））、拍摄条件敏感性、以及动作类型敏感性。综合三类技术方案的比较表明，LPT以直接位移测量原理占据精度最高层级，计算机视觉方案在标准化条件下可接近LPT精度，但在真实场景中受拍摄条件与漏检率影响，稳定性尚待提升。

\section{主观反馈路径的方法学演进与实证进展}

主观反馈路径通过训练者自身的内在评估来调节负荷，不需要外部测量设备。其代表方法经历了从RPE到RIR的演进，并因ACSM 2026立场声明的推荐而获得权威认可。

\subsection{从RPE到RIR的方法学演进}

\citep{Borg1982}提出的RPE是主观反馈路径的开创性工具。\citep{Zourdos2016}针对抗阻训练的动作特点开发了基于RIR的量表，将抽象的用力感转化为"还能完成几次"这一有直接操作价值的判断。RIR相较于传统RPE的方法学优势体现在操作直接性、与动作力学的直接对应性以及与客观参数的可桥接性三个方面。\citep{Jukic2024}通过对RIR与杠铃速度关系的个体化建模证实，个体化的RIR-速度关系可实现可接受的RIR预测准确性，这为主观反馈与客观反馈之间的方法学融合提供了实证基础。

\subsection{RIR训练的效度证据与力竭训练争议}

\citep{Grgic2022}的Meta分析证据表明，训练至力竭并非获得力量、肌肥大或功率适应的必要条件；保留一定RIR的非力竭训练在力量与肌肥大适应上与力竭训练无显著差异。\citep{Refalo2024}的随机对照试验进一步显示，RIR组在获得等效适应的同时，具有更低的主观疲劳与更高的依从性。基于上述证据，当前研究普遍认为大多数训练情境下采用保留1至3 RIR的非力竭训练是合理选择。

然而，RIR判断的准确性受训练经验影响较大。\citep{Hackett2017}指出，训练经验较少者的RIR预测误差可达2至3次；而训练经验丰富者的预测误差通常可控制在1次以内。这一差异说明，RIR的效度并非固定不变，而是随训练者经验水平而变化。

\subsection{ACSM 2026立场声明的方法学转向}

\citep{Currier2026}于2026年4月发布的最新抗阻训练立场声明距上一版本发布已有17年，是ACSM对抗阻训练处方框架的一次系统性更新。立场声明最核心的方法学转向体现在强度处方指标的更换上：明确推荐以RIR取代\%1RM作为抗阻训练强度处方的首选指标。同时，\citep{Currier2026}明确推荐大多数训练情境下采用保留1至3 RIR的非力竭训练，并在立场声明层面明确提出"训练接近但不必达到力竭"的处方立场。

\citep{Currier2026}明确将其推荐对象限定为"经验丰富的健康成年抗阻训练者"，其证据基础主要来自具备一定训练经验的健康人群。在训练初学者、老年人群与临床康复人群中的适用性，仍需谨慎评估。

\section{两条反馈路径的直接对比与实际应用权衡}

两条路径在各自参照框架内都已建立较充分的证据基础：VBT相较于PBT的适应性优势获得多层级证据支持，RIR训练与力竭训练的等效性也获得Meta分析和权威指南的一致确认。但两条路径之间的直接对比证据相对稀缺，且现有证据呈现出与各自独立证据体系不完全一致的复杂图景。

\subsection{直接对比研究的证据现状}

\citep{Naclerio2025}率先在卧推动作中比较了基于速度反馈（VEL）的客观自主调节方案与基于RPE的主观自主调节方案，结果显示两种方案在总训练量、平均速度与速度损失控制等核心训练产出指标上均未表现出显著差异，研究者将这一结果概括为"可比的训练产出"。\citep{Naclerio2026}随后将相同设计拓展至深蹲动作，重复了卧推研究的核心发现。

\citep{Orange2022}在其关于VBT与传统训练方法对比的系统综述中专门分析了纳入研究中涉及主观反馈对照组的数据，指出当对照组采用RPE或RIR等主观调节方案时，VBT组的适应性优势显著缩小甚至消失。2025年发表的一项涵盖APRE、RPE、VBRT与PBT四种自主调节方案的网络Meta分析进一步量化了这一现象，结果显示VBRT与RPE之间的效应量差异仅为SMD=0.28，属于小效应。

\subsection{可比结果的机制解释}

第一种可能的解释是，两条路径最终都在控制力竭接近度。无论是VBT通过速度损失阈值控制组内疲劳，还是RIR通过主观感知控制距力竭距离，两条路径最终都在执行同一项核心功能：将训练者的力竭接近度维持在特定目标区间内。

第二种可能的解释涉及干预周期与测量粒度的限制。现有直接对比研究的干预周期普遍较短，Naclerio系列仅为2周急性训练产出评估，可能不足以放大两条路径在长期适应上的潜在差异。

第三种可能的解释涉及反馈的行为学效应差异。\citep{Wulf2016}的OPTIMAL理论指出，客观速度反馈通过增强期望与外部注意焦点优化运动表现，而主观反馈通过自主性支持增强内在动机。两条路径在行为学层面的作用方式存在差异，但这些差异可能不会直接转化为短期力量或爆发力适应的组间差异。

\subsection{既有研究的主要局限}

综合上述分析，既有文献在以下四个方面存在明显局限。第一，VBT与规范化RIR训练的直接对比证据极为稀缺。第二，"可比"结果的解释尚未建立。第三，真实训练场景的生态效度证据不足。第四，训练者主观体验维度的系统性缺失。

本研究的设计直接回应上述四个局限。研究二采用两组平行随机对照预试验设计，以基于计算机视觉AI系统的客观速度反馈组与基于RIR的规范化主观反馈组作为直接对比条件；研究一通过移动端AI系统与专业LPT设备的效度验证，系统评估低成本CV方案在真实训练场景中的测量可行性；第三阶段通过半结构化深度访谈与主题分析，从主观体验层面揭示训练者在两条反馈路径下的决策逻辑与工具接受度。
